Let $(X,d)$ with $X=\{P_1,\dots,P_n\}$ be a finite metric space. We will call each pair $(P_i,P_j)$ $(i\ne j)$ an edge and the distance between them $d(P_i,P_j)$ the {\em edge length} and denote it by $d_{ij}$. 
Let $tX$ $(t>0)$ denote the metric space $(X,td)$. 
The {\em similarity matrix} $Z_X(t)$ is given by $Z_X(t)=\left(\exp(-t\, d_{ij})\right)_{i,j}$. %
In this paper we only use the cases when $Z_X(t)$ is invertible. 
Then the {\em magnitude function}, denoted by $|tX|$ or $M_X(t)$, is given by the sum of all the entries of $Z_X(t)^{-1}$. 
It is known that $M_X(t)$ is an increasing function of $t$ for $t\gg0$ and that $\lim_{t\to\infty}M_X(t)=\#X$ (\cite{L13} Proposition 2.2.6). 
A space $X$ is said to have {\em one-point property} if $\lim_{t\to0^+}M_X(t)=1$. 
Any compact subset of Euclidean space has one-point property (\cite{LM23} Theorem 3.1). 

An alternative expression of the magnitude can be obtained by putting $q=e^{-t}$ (\cite{L19}). The similarity matrix is then given by $z_X(q)=\left(q^{d_{ij}}\right)_{i,j}$. 
The determinant and the sum of all the cofactors of $z_X(q)$ are both ``{\em generalized polynomials}'' that allow non-integer exponents. 
Since the determinant has the constant term $1$, it is invertible in the field of generalized rational functions. The sum of all the entries of $z_X(q)^{-1}$ is called the {\em formal magnitude}, and denoted by $m_X(q)$. 
In this article we assume that $m_X(q)$ is expanded as a ``{\em generalized formal power series}'' that allows non-integer exponents. 
